\documentclass[12pt]{article}
\usepackage{amsmath,amssymb,amsthm,hyperref,booktabs,xcolor,enumitem}
\usepackage[margin=1in]{geometry}

\newtheorem{theorem}{Theorem}
\newtheorem{lemma}[theorem]{Lemma}
\newtheorem{definition}[theorem]{Definition}
\newtheorem{remark}[theorem]{Remark}
\newtheorem{proposition}[theorem]{Proposition}
\newtheorem{corollary}[theorem]{Corollary}

\definecolor{proved}{RGB}{40,120,40}
\definecolor{open}{RGB}{160,50,50}
\definecolor{partial}{RGB}{160,120,0}

\title{\textbf{Exploration: \texttt{sin\_not\_in\_real\_EML\_k}}\\[0.4em]
\large Initial Attack: Structural and Analytic Obstructions\\
\normalsize (Second Remaining Sorry)}
\author{Arturo R.\ Almaguer}
\date{2026-04-21}

\begin{document}
\maketitle

\begin{abstract}
We explore the second remaining sorry in \texttt{InfiniteZerosBarrier.lean}:
\texttt{sin\_not\_in\_real\_EML\_k}, which states that the function
$z \mapsto \sin(z.\mathtt{re})$ is not in $\mathrm{EML}_k$ (the set of complex
EML tree functions of depth $\leq k$) for any finite $k$.
We map the relationship between this sorry and \texttt{sin\_not\_in\_eml},
identify early structural obstructions provable without o-minimal theory,
and characterize exactly where o-minimality enters.
\end{abstract}

%──────────────────────────────────────────────────────────────────────────────
\section{The Target Sorry}

\begin{verbatim}
theorem sin_not_in_real_EML_k (k : N) :
    (fun x : C => (Real.sin x.re : C)) NOT in EML_k k := by
  sorry
\end{verbatim}

Here $\mathrm{EML}_k$ is the set of $\mathbb{C} \to \mathbb{C}$ functions:
\[
  \mathrm{EML}_k = \{f : \mathbb{C} \to \mathbb{C} \mid \exists t : \mathtt{EMLTree},\ t.\mathtt{depth} \leq k \land \forall x,\ f(x) = t.\mathtt{eval}(x)\}.
\]
The target function is $F(z) = \sin(z.\mathtt{re})$ — real sine applied to the real part of $z$.

%──────────────────────────────────────────────────────────────────────────────
\section{Relation to sin\_not\_in\_eml}

\begin{lemma}[Reduction]\label{lem:reduction}
\texttt{sin\_not\_in\_eml} implies \texttt{sin\_not\_in\_real\_EML\_k}.
\end{lemma}

\begin{proof}
Suppose $F \in \mathrm{EML}_k$, i.e., there exists $t$ with $t.\mathtt{depth} \leq k$
and $F(z) = t.\mathtt{eval}(z)$ for all $z \in \mathbb{C}$.
Restrict to real inputs: for $x \in \mathbb{R}$, $F(\uparrow\!x) = \sin(x)$.
And $t.\mathtt{eval}(\uparrow\!x) = \uparrow(t.\mathtt{evalReal}(x))$ when the
tree is real-valued on real inputs (\textbf{see gap below}).
If we can equate the real parts, $t.\mathtt{evalReal}(x) = \sin(x)$ for all $x \in \mathbb{R}$.
This contradicts \texttt{sin\_not\_in\_eml}. \qed (modulo the gap)
\end{proof}

\subsection{The Reduction Gap}

The step ``$t.\mathtt{eval}(\uparrow\!x) = \uparrow(t.\mathtt{evalReal}(x))$'' requires
that on real inputs, $t$ produces real outputs. This holds when $t$ has real constants
and real log arguments everywhere. It is not automatic for trees with complex constants.

Specifically, $t.\mathtt{evalReal}(x) = (t.\mathtt{eval}(\uparrow\!x)).\mathtt{re}$
by definition. So $t.\mathtt{eval}(\uparrow\!x) = \uparrow(t.\mathtt{evalReal}(x))$
requires $(t.\mathtt{eval}(\uparrow\!x)).\mathtt{im} = 0$.

\begin{remark}[When real-valuedness holds]
If all constants in $t$ are real ($c \in \mathbb{R}$) and all log arguments are
positive reals on $\mathbb{R}_{>0}$, then by induction, $t.\mathtt{eval}(\uparrow\!x)$
is real for all $x > 0$. This follows from:
\begin{itemize}[noitemsep]
  \item $\mathtt{const}\ c$: $c \in \mathbb{R}$ means $c.\mathtt{im} = 0$.
  \item $\mathtt{var}$: $\uparrow\!x$ is real.
  \item $\mathtt{ceml}\ t_1\ t_2$: if $t_1, t_2$ produce reals and $t_2 > 0$,
    then $\exp(\text{real}) - \ln(\text{positive real})$ is real.
\end{itemize}
\end{remark}

%──────────────────────────────────────────────────────────────────────────────
\section{Structural Obstructions Provable Now}

\subsection{Obstruction 1: $F$ is not complex-analytic}

\begin{proposition}[\textcolor{proved}{Provable now}]\label{prop:not-analytic}
$F(z) = \sin(z.\mathtt{re})$ is \emph{not} complex-analytic at any point.
\end{proposition}

\begin{proof}
$F(x + iy) = \sin(x)$ depends only on the real part. By the Cauchy-Riemann equations,
if $F$ were analytic, $\partial_x \sin(x) = \cos(x)$ would equal $\partial_y 0 = 0$,
requiring $\cos(x) = 0$ for all $x$ — contradiction. \qed
\end{proof}

\begin{corollary}[\textcolor{proved}{Provable now}]\label{cor:not-holomorphic}
If any $t \in \mathrm{EML}_k$ computes only entire (holomorphic everywhere) functions,
then $F \notin \mathrm{EML}_k$ immediately.
\end{corollary}

\textbf{Key question}: Are all EML tree functions entire?
\begin{itemize}[noitemsep]
  \item \texttt{const}, \texttt{var}: entire.
  \item $\mathtt{ceml}\ t_1\ t_2$: $\exp(t_1.\mathtt{eval})$ is entire (comp of entire).
    $\log(t_2.\mathtt{eval})$ is analytic only where $t_2 \neq 0$ and $t_2 \notin \mathbb{R}_{\leq 0}$.
\end{itemize}

So EML trees are NOT entire in general (log introduces branch cuts).
The Corollary \ref{cor:not-holomorphic} does not immediately apply.

\subsection{Obstruction 2: Restriction to the real line}

\begin{proposition}[\textcolor{proved}{Provable now — modulo depth-1 result}]
$F \notin \mathrm{EML}_1$.
\end{proposition}

\begin{proof}
If $t \in \mathrm{EML}_1$ with $t.\mathtt{eval} = F$ everywhere on $\mathbb{C}$,
then in particular $t.\mathtt{eval}(\uparrow\!x) = F(\uparrow\!x) = \uparrow(\sin(x))$
for all real $x$. Taking real parts: $t.\mathtt{evalReal}(x) = \sin(x)$.
But $t$ has depth $\leq 1$, and $\sin \notin \mathrm{EML}_1(\mathbb{R})$
(Proposition 2 of \texttt{sin\_not\_in\_eml\_o\_minimal\_Start.tex}).
Contradiction. \qed (modulo the real-part extraction step for depth-1 complex trees)
\end{proof}

\subsection{Obstruction 3: Infinite zeros on the real slice}

\begin{proposition}[\textcolor{partial}{Partial — needs general zero-count bound}]
For any $t$ with $t.\mathtt{eval} = F$:
$\{x \in \mathbb{R} \mid t.\mathtt{evalReal}(x) = 0\}$ is infinite.
\end{proposition}

\begin{proof}
$t.\mathtt{evalReal}(x) = F(\uparrow\!x).\mathtt{re} = \sin(x).\mathtt{re} = \sin(x)$.
$\sin(n\pi) = 0$ for all $n \in \mathbb{Z}$, giving infinitely many real zeros. \qed
\end{proof}

If we additionally prove that any EML tree has finitely many zeros on $\mathbb{R}$
(the depth-$k$ version of \texttt{depth1\_finite\_zeros\_real}), contradiction follows.
That depth-$k$ version is the same o-minimal gap as in \texttt{sin\_not\_in\_eml}.

%──────────────────────────────────────────────────────────────────────────────
\section{The Core Common Difficulty}

Both remaining sorries reduce to the same statement:

\begin{center}
\fbox{\begin{minipage}{0.85\linewidth}
\textbf{Key Lemma} (not yet proved for general $k$):
For any $t : \mathtt{EMLTree}$ with $t.\mathtt{depth} \leq k$,
the zero set of $t.\mathtt{evalReal}$ is finite on any bounded interval.
\end{minipage}}
\end{center}

For depth $\leq 1$: \textcolor{proved}{proved} (\texttt{depth1\_finite\_zeros\_real}, 0 sorry).
For general $k$: \textcolor{open}{open} (needs o-minimal theory or depth-by-depth analysis).

%──────────────────────────────────────────────────────────────────────────────
\section{What Changes Between the Two Sorries}

\begin{center}
\begin{tabular}{lll}
\toprule
Aspect & \texttt{sin\_not\_in\_eml} & \texttt{sin\_not\_in\_real\_EML\_k} \\
\midrule
Domain of $t$ & $\mathbb{R} \to \mathbb{R}$ (real) & $\mathbb{C} \to \mathbb{C}$ (complex) \\
Target function & $\sin : \mathbb{R} \to \mathbb{R}$ & $z \mapsto \sin(z.\mathtt{re})$ \\
Proof by reduction & No (direct) & Yes (restrict to $\mathbb{R}$) \\
Additional gap & WFP derivation & Real-part extraction \\
Core difficulty & Same & Same \\
\bottomrule
\end{tabular}
\end{center}

\texttt{sin\_not\_in\_real\_EML\_k} is logically weaker: it follows from
\texttt{sin\_not\_in\_eml} once we establish that for trees equaling $F$ on $\mathbb{C}$,
the real restriction equals $\sin$ on $\mathbb{R}$. This restriction step is cleaner
for the special case of trees with real constants.

%──────────────────────────────────────────────────────────────────────────────
\section{Near-Term Action Plan}

\begin{enumerate}
  \item \textbf{Prove $F \notin \mathrm{EML}_1$ in Lean} (provable now):
    Use depth-1 real restriction + \texttt{depth1\_finite\_zeros\_real}.

  \item \textbf{Prove $F \notin \mathrm{EML}_2$ by case analysis}:
    Enumerate all depth-$\leq 2$ tree shapes. Each produces a function with
    finitely many real zeros; $\sin$ does not. More cases but still formalizable.

  \item \textbf{Factor out the common Key Lemma}:
    State \texttt{EML\_depth\_k\_finite\_zeros} as a standalone sorry
    that, once proved, gives both \texttt{sin\_not\_in\_eml} and
    \texttt{sin\_not\_in\_real\_EML\_k} immediately.

  \item \textbf{Long-term}: Formalize Khovanskii bound or Wilkie's theorem.
\end{enumerate}

%──────────────────────────────────────────────────────────────────────────────
\section{Summary}

\begin{center}
\begin{tabular}{lll}
\toprule
Claim & Status & Path \\
\midrule
$F \notin \mathrm{EML}_1$ & \textcolor{proved}{Provable now} & Real restriction + depth-1 \\
$F \notin \mathrm{EML}_2$ & \textcolor{partial}{Feasible} & Case analysis \\
$F \notin \mathrm{EML}_k$ (all $k$) & \textcolor{open}{Open} & O-minimal theory \\
$F$ not complex-analytic & \textcolor{proved}{Proved} & Cauchy-Riemann \\
Real zeros of $t$ equal real zeros of $\sin$ & \textcolor{proved}{Provable} & Real restriction \\
\bottomrule
\end{tabular}
\end{center}

The sorry is not an engineering gap. It is a genuine mathematical gap in the same
class as \texttt{sin\_not\_in\_eml}, sharing the same core lemma about zero-count
bounds. Both will be resolved together once the depth-$k$ zero-finiteness result
is established.

\end{document}
